\documentclass{article}

\usepackage{amsmath,amssymb}
\usepackage{xurl}
\usepackage[hidelinks]{hyperref}
\setlength{\emergencystretch}{2em}

\input{command}

\title{Wolf Proposition~6.1 --- The Russo--Dye Factor}
\author{}
\date{2026-08-23}

\begin{document}
\maketitle
\gapnote{scope-restriction}{resolved}

\section{The assertion}

Let $d\geq1$, and let $T:\MN{d}\to\MN{d}$ be positive.
Proposition~6.1 of Ref.~\cite{Wolf2012Quantum} states
\begin{align}
    \varrho(T)
        &\leq\|T(\Id)\|_\infty.
    \label{eq:wolf_prop61_russo_dye_factor_spectral_radius_bound}
\end{align}
This is Equation~(6.2) of the source. If $T$ is unital or trace preserving,
then $1$ is an eigenvalue, $\varrho(T)=1$, and every eigenvalue lies in the
closed unit disk.\footnote{Local source:
\path{Notes/WolfNoteTexSource/ch06_spectral_properties.tex}, lines~73--91.}

The proof invokes the Russo--Dye estimate
\begin{align}
    \|T(X)\|_\infty
        &\leq\|T(\Id)\|_\infty\,\|X\|_\infty.
    \label{eq:wolf_prop61_russo_dye_factor_russo_dye_estimate}
\end{align}
This occurs at line~84 of the local source. If $X\neq0$ and
\begin{align}
    T(X)
        &=\lambda X,
    \label{eq:wolf_prop61_russo_dye_factor_eigenvector_equation}
\end{align}
then
\begin{align}
    |\lambda| \|X\|_\infty
        &=\|T(X)\|_\infty,
    \label{eq:wolf_prop61_russo_dye_factor_eigenvalue_norm_identity}\\
    \|T(X)\|_\infty
        &\leq\|T(\Id)\|_\infty\,\|X\|_\infty,
    \label{eq:wolf_prop61_russo_dye_factor_eigenvalue_norm_estimate}\\
    |\lambda|
        &\leq\|T(\Id)\|_\infty.
    \label{eq:wolf_prop61_russo_dye_factor_eigenvalue_bound}
\end{align}
The final inequality is Equation~(6.3) of
Ref.~\cite{Wolf2012Quantum}. Maximizing over the eigenvalue moduli gives
(\ref{eq:wolf_prop61_russo_dye_factor_spectral_radius_bound}).

\section{The formalization}

The sharp coefficient-one theorem is formalized in
\doclink{QICLean/Channel/Peripheral/SpectralRadius}. Its principal
declarations are as follows.
\begin{itemize}
    \item
    \leanid{IsPositiveMap.norm\_apply\_le\_norm\_map\_one\_mul\_norm}
    proves
    (\ref{eq:wolf_prop61_russo_dye_factor_russo_dye_estimate}) for every
    matrix $X$.

    \item
    \leanid{IsPositiveMap.eigenvalue\_norm\_le\_norm\_map\_one}
    proves
    (\ref{eq:wolf_prop61_russo_dye_factor_eigenvalue_bound}).

    \item
    \leanid{IsPositiveMap.spectralRadius\_le\_nnnorm\_map\_one}
    proves
    (\ref{eq:wolf_prop61_russo_dye_factor_spectral_radius_bound}).

    \item
    \leanid{eigenvalue\_one\_of\_map\_one\_eq\_one},
    \leanid{IsPositiveMap.eigenvalue\_norm\_le\_one\_of\_map\_one\_eq\_one},
    and
    \leanid{IsPositiveMap.spectralRadius\_eq\_one\_of\_map\_one\_eq\_one}
    prove the three unital conclusions without assuming complete positivity.
\end{itemize}

The trace-preserving results remain independent sibling theorems. The
declaration
\leanid{IsPositiveMap.eigenvalue\_norm\_le\_one\_of\_tracePreserving} in
\doclink{QICLean/Channel/Determinant/Bound} gives the closed-unit-disk bound,
and
\leanid{IsPositiveMap.eigenvalue\_one\_exists\_of\_tracePreserving} produces
a nonzero positive-semidefinite fixed point. The declaration
\leanid{IsPositiveMap.spectralRadius\_eq\_one\_of\_tracePreserving} combines
these results to prove the trace-preserving spectral-radius equality.

\section{The proof route}

The pinned Mathlib version contains neither a standalone Russo--Dye theorem
nor a theorem identifying the operator-norm unit ball with the closed convex
hull of the unitaries. The formal proof instead derives
(\ref{eq:wolf_prop61_russo_dye_factor_russo_dye_estimate}) from
Theorem~5.6 of Ref.~\cite{Wolf2012Quantum}, formalized as
\leanid{KadisonSchwarz.schwarz\_inequality\_commuting\_dominant\_operator}.
This is an alternative derivation of the estimate used in the source, not a
formalization of the convex-hull theorem.

Set
\begin{align}
    c
        &:=\|T(\Id)\|_\infty.
    \label{eq:wolf_prop61_russo_dye_factor_normalization_constant}
\end{align}
For $c>0$, define
\begin{align}
    S
        &:=c^{-1}T.
    \label{eq:wolf_prop61_russo_dye_factor_normalized_map}
\end{align}
The map $S$ is positive and subunital. If $\|A\|_\infty\leq1$, then
\begin{align}
    A^\dagger A
        &\leq\Id.
    \label{eq:wolf_prop61_russo_dye_factor_unit_ball_order}
\end{align}
Applying Theorem~5.6 of Ref.~\cite{Wolf2012Quantum} with dominant operator
$\Id$ gives
\begin{align}
    S(A)^\dagger S(A)
        &\leq S(\Id),
    \label{eq:wolf_prop61_russo_dye_factor_schwarz_first_bound}\\
    S(\Id)
        &\leq\Id,
    \label{eq:wolf_prop61_russo_dye_factor_subunital_bound}\\
    S(A)^\dagger S(A)
        &\leq\Id.
    \label{eq:wolf_prop61_russo_dye_factor_schwarz_unit_bound}
\end{align}
The C$^*$-identity then yields
\begin{align}
    \|S(A)\|_\infty
        &\leq1.
    \label{eq:wolf_prop61_russo_dye_factor_normalized_contraction}
\end{align}

The case $c=0$ is separate. The same inequality gives
\begin{align}
    T(A)^\dagger T(A)
        &\leq T(\Id),
    \label{eq:wolf_prop61_russo_dye_factor_zero_case_schwarz}\\
    T(\Id)
        &=0,
    \label{eq:wolf_prop61_russo_dye_factor_zero_case_identity}\\
    T(A)^\dagger T(A)
        &\leq0,
    \label{eq:wolf_prop61_russo_dye_factor_zero_case_nonpositive_square}\\
    T(A)
        &=0.
    \label{eq:wolf_prop61_russo_dye_factor_zero_case_map_value}
\end{align}
Rescaling an arbitrary nonzero $X$ then proves
(\ref{eq:wolf_prop61_russo_dye_factor_russo_dye_estimate}) in all cases.
The argument does not decompose $X$ into four positive parts and introduces
no factor four.

The formal hypotheses match the source: $T$ is positive and $d\geq1$. The
norm is Mathlib's matrix C$^*$-operator norm, corresponding to
$\|\cdot\|_\infty$ in Ref.~\cite{Wolf2012Quantum}; it is neither the
Frobenius norm nor the row-sum matrix norm.

\section{Verdict}

The former scope restriction is resolved. Equations~(6.2) and~(6.3), and the
unital conclusions of Proposition~6.1 in
Ref.~\cite{Wolf2012Quantum}, are formalized for arbitrary positive matrix
maps with sharp coefficient one and without a complete-positivity assumption.
The formal proof derives the source's norm estimate from Theorem~5.6 of the
same reference.\footnote{Tracking: \ghissue{36}.}

\bibliographystyle{alpha}
\bibliography{references}

\end{document}
